\documentclass[11pt,a4paper,reqno]{amsart}

\usepackage[utf8]{inputenc}
\usepackage[T1]{fontenc}
\usepackage{amsmath,amssymb,amsfonts,amsthm,mathtools}
\usepackage{mathrsfs}
\usepackage{microtype}
\usepackage{booktabs}
\usepackage{tabularx}
\usepackage[dvipsnames]{xcolor}
\usepackage{hyperref}
\usepackage{geometry}
\usepackage{enumitem}
\usepackage{cite}

\geometry{
  top=3cm,
  bottom=3cm,
  left=2.8cm,
  right=2.8cm,
  headheight=14pt
}

\hypersetup{
  colorlinks=true,
  linkcolor=NavyBlue,
  citecolor=Mahogany,
  urlcolor=TealBlue,
  pdfauthor={Reinaldo M. Silva-Filho},
  pdftitle={Geometric Flows, Partial Differential Equations, and Variational Dynamics on Matrix and Tensor Manifolds}
}

% --- Theorem Environments ---
\newtheorem{theorem}{Theorem}[section]
\newtheorem{lemma}[theorem]{Lemma}
\newtheorem{proposition}[theorem]{Proposition}
\newtheorem{corollary}[theorem]{Corollary}

\theoremstyle{definition}
\newtheorem{definition}[theorem]{Definition}
\newtheorem{example}[theorem]{Example}
\newtheorem{remark}[theorem]{Remark}
\newtheorem{construction}[theorem]{Construction}

\numberwithin{equation}{section}

% --- Custom Notations ---
\providecommand{\doi}[1]{\href{https://doi.org/#1}{\texttt{doi:#1}}}
\providecommand{\R}{\mathbb{R}}
\newcommand{\C}{\mathbb{C}}
\newcommand{\N}{\mathbb{N}}
\newcommand{\Z}{\mathbb{Z}}
\providecommand{\Sph}{\mathbb{S}}
\newcommand{\Tor}{\mathbb{T}}
\newcommand{\Id}{\mathrm{Id}}
\newcommand{\Tr}{\operatorname{Tr}}
\newcommand{\rank}{\operatorname{rank}}
\newcommand{\supp}{\operatorname{supp}}
\newcommand{\Lip}{\operatorname{Lip}}
\providecommand{\BV}{\operatorname{BV}}
\providecommand{\TV}{\operatorname{TV}}
\newcommand{\Per}{\operatorname{Per}}
\newcommand{\Ric}{\operatorname{Ric}}
\providecommand{\cutnorm}[1]{\|#1\|_{\square}}
\newcommand{\norm}[1]{\left\|#1\right\|}
\newcommand{\inner}[2]{\left\langle #1, #2 \right\rangle}
\newcommand{\dif}{\,\mathrm{d}}
\newcommand{\abs}[1]{\left|#1\right|}

\title[Geometric Flows on Matrix and Tensor Manifolds]{Geometric Flows, Partial Differential Equations, and Variational Dynamics on Matrix and Tensor Manifolds}

\author{Reinaldo M. Silva-Filho}
\address{Universidade Federal de Lavras (UFLA), Departamento de Estat\'istica e Experimenta\c{c}\~ao Agropecu\'aria (PPGEEAA/DES)}
\email{reinaldo.filho1@estudante.ufla.br}
\date{\today}


\DeclareMathOperator*{\argmin}{argmin}
\providecommand{\Hyp}{\mathbb{H}}
\providecommand{\Sph}{\mathbb{S}}
\providecommand{\II}{\mathrm{II}}
\providecommand{\R}{\mathbb{R}}
\providecommand{\Area}{\mathrm{Area}}
\providecommand{\Hsn}{\mathcal{H}}
\providecommand{\BV}{\mathrm{BV}}
\providecommand{\TV}{\mathrm{TV}}
\providecommand{\cutnorm}[1]{\|#1\|_{\square}}

\begin{document}

\begin{abstract}
While static functional realizations map finite discrete arrays into infinite-dimensional spaces, dynamical processes in high-dimensional data, network physics, and optimization require a non-equilibrium geometric foundation. In this treatise, we establish a rigorous framework for \emph{Geometric Flows, Non-Local Partial Differential Equations, and Variational Dynamics} operating on matrix and tensor manifolds, as well as their infinite-dimensional continuum limits. 

We systematically investigate: (i) canonical Riemannian metrics on matrix manifolds---including the affine-invariant geometry of the positive definite cone $\mathcal{S}_{++}^n$, Bures-Wasserstein optimal transport geometries, and the sub-Riemannian structure of fixed-rank matrix varieties $\mathcal{M}_r$; (ii) continuous Hamiltonian and gradient dynamics, demonstrating that the celebrated Toda lattice flow interpolates the discrete QR algorithm via isospectral Lax-pair equations; (iii) the continuum theory of \emph{Graphon Partial Differential Equations}, where dense networks evolving in time converge to non-local diffusion equations governed by continuous graphon Laplacians $\mathcal{L}_W$, rigorously establishing parabolic smoothing, dissipation of cut norms, and contraction of graph entropy; and (iv) curvature-driven geometric evolutions, constructing the continuous \emph{Ollivier-Wasserstein Ricci Flow} on graphon moduli spaces and coupling the geometric Coarea formula to Mean Curvature Flows on matrix level sets. Finally, we formulate the continuous renormalization group flow of Tensor Network states (MPS/PEPS) as an energy-dissipating Riemannian gradient flow, providing a unified geometric bridge between continuous PDEs and discrete multilinear algorithms.
\end{abstract}

\keywords{Matrix manifolds, geometric flows, Riemannian optimization, Toda lattice, graphon PDEs, non-local diffusion, Ricci flow, Ollivier curvature, Tensor Networks.}
\subjclass[2020]{53C44, 35K90, 15A69, 47J35, 49Q20, 58B20}

\maketitle

\tableofcontents

\newpage

% =========================================================================
\section{Introduction: From Static Realizations to Geometric Dynamics}
\label{sec:intro}
% =========================================================================

In the companion work \cite{silvafilho2026functional}, we investigated the static embedding of discrete multilinear arrays into function spaces via Functional Realization Operators $\Phi \colon \mathcal{T}^k(V) \to \mathcal{F}(\Omega)$, revealing topological and geometric invariants (Sobolev energies, total variation, Morse spectra, and graphon cut limits) inaccessible to pure linear algebra. 

However, in computational mathematics, statistical physics, and non-convex optimization, matrices and tensors rarely remain static. Rather, they evolve continuously in time $t \ge 0$ under continuous dynamical laws:
\begin{equation}
    \label{eq:matrix_ode_general}
    \frac{\dif A(t)}{\dif t} = \mathcal{V}(A(t)), \quad A(0) = A_0,
\end{equation}
where $\mathcal{V}$ is a vector field on an ambient matrix space or tangent to a curved sub-manifold $\mathcal{M}$. Classical examples include continuous-time neural learning, natural gradient descent in machine learning, quantum state evolution via Lindbladian master equations, and the continuous interpolation of numerical linear algebra algorithms such as the QR algorithm and power iterations.

Despite the ubiquity of these dynamical processes, an overarching geometric and PDE-theoretic formulation unifying discrete matrix flows with continuum partial differential equations remains largely fragmented. Specifically:
\begin{enumerate}[label=(\alph*)]
    \item \textbf{Curvature and Invariance on Matrix Cones:} The Euclidean Frobenius metric $\inner{U}{V}_F = \Tr(U^T V)$ is flat, yet the space of positive definite covariance matrices $\mathcal{S}_{++}^n$ inherently possesses negative Riemannian sectional curvature under the affine-invariant metric, profoundly shaping convergence rates of gradient flows.
    \item \textbf{Rank Degeneracy and Boundary Singularities:} The variety of real $m \times n$ matrices of fixed rank $r < \min(m, n)$, denoted $\mathcal{M}_r$, is an embedded Riemannian submanifold. However, its closure $\overline{\mathcal{M}}_r$ contains matrices of rank $< r$ as singular boundary strata, leading to curvature blow-ups and loss of tangent projection well-posedness in matrix completion algorithms.
    \item \textbf{The Continuum Limit of Evolving Networks ($n \to \infty$):} While graphons $W \colon [0, 1]^2 \to [0, 1]$ characterize static dense graph limits in the cut metric, dynamic processes on large networks require solving non-local partial differential equations $\partial_t W(t, x, y) = \mathcal{F}(W(t, \cdot, \cdot))$.
    \item \textbf{Curvature-Driven Matrix Smoothing:} While Ricci flows $\partial_t g = -2 \Ric(g)$ deform Riemannian metrics to uniformize geometry, their matrix and graphon analogues---driven by Ollivier-Wasserstein coarse curvature---remain to be systematized.
\end{enumerate}

\subsection*{Core Contributions}
In this treatise, we develop the differential geometry and PDE theory of matrix and tensor flows. We prove:
\begin{enumerate}[label=\arabic*.]
    \item The equivalence of Toda lattice Hamiltonian flows with the continuous QR eigenvalue algorithm under isospectral commutator dynamics (Theorem~\ref{thm:toda_qr});
    \item Strict rank preservation and energy decay for projected Riemannian gradient flows on $\mathcal{M}_r$ prior to reaching the boundary variety $\overline{\mathcal{M}}_{r-1}$ (Theorem~\ref{thm:rank_preservation});
    \item Global well-posedness, contraction of the cut norm, and singular support decay for non-local Graphon Heat Equations $\partial_t W = \Delta_\otimes W$ (Theorems~\ref{thm:graphon_heat_existence} and \ref{thm:heat_smoothing});
    \item The asymptotic neckpinch disconnection in graphon Ricci flows driven by Ollivier curvature concentration (Theorem~\ref{thm:ricci_neckpinch});
    \item The variational convergence of matrix level-set boundaries under discrete mean curvature flows via the Coarea identity (Theorem~\ref{thm:mean_curvature_coarea}).
\end{enumerate}

% =========================================================================
\section{Riemannian Metrics and Geometry on Matrix and Tensor Manifolds}
\label{sec:riemannian_geometry}
% =========================================================================

Let $\R^{m \times n}$ denote the space of real matrices. We endow matrix spaces and subsets thereof with non-Euclidean Riemannian structures.

\subsection{The Cone of Symmetric Positive Definite Matrices}
\label{subsec:cone_spd}
The set $\mathcal{S}_{++}^n \coloneqq \{A \in \operatorname{Sym}_n(\R) \colon A \succ 0\}$ is an open convex cone in $\R^{n \times n}$. While the Euclidean metric yields flat zero curvature, the natural geometry arises from the transitive Lie group action of $GL(n)$ via congruence $A \mapsto G A G^T$.

\begin{definition}[Affine-Invariant Metric]
For $A \in \mathcal{S}_{++}^n$ and tangent vectors $U, V \in T_A \mathcal{S}_{++}^n \cong \operatorname{Sym}_n(\R)$, the \emph{Affine-Invariant Riemannian Metric} is defined by:
\begin{equation}
    \label{eq:affine_metric}
    g_A^{\mathrm{AI}}(U, V) \coloneqq \Tr\left( A^{-1} U A^{-1} V \right) = \Tr\left( A^{-1/2} U A^{-1} V A^{-1/2} \right).
\end{equation}
\end{definition}

\begin{theorem}[Curvature of the Affine-Invariant Cone]
\label{thm:cone_curvature}
Endowed with the affine-invariant metric $g^{\mathrm{AI}}$, the space $(\mathcal{S}_{++}^n, g^{\mathrm{AI}})$ is a complete, simply connected Hadamard manifold (symmetric space $GL(n)/O(n)$) with strictly non-positive Riemannian sectional curvature:
\begin{equation}
    K(U, V) \le 0, \quad \forall U, V \in T_A \mathcal{S}_{++}^n.
\end{equation}
The unique geodesic connecting $A, B \in \mathcal{S}_{++}^n$ is given explicitly by:
\begin{equation}
    \label{eq:cone_geodesic}
    \gamma(t) = A^{1/2} \left( A^{-1/2} B A^{-1/2} \right)^t A^{1/2}, \quad t \in [0, 1].
\end{equation}
\end{theorem}
\begin{proof}
Let $A(t)$ be a curve in $\mathcal{S}_{++}^n$ with $A(0) = A$ and $\dot{A}(0) = U$. The Levi-Civita connection $\nabla$ associated with \eqref{eq:affine_metric} is given by:
\begin{equation}
    \nabla_U V = D V[U] - \frac{1}{2}(U A^{-1} V + V A^{-1} U),
\end{equation}
where $D V[U]$ denotes the directional Euclidean derivative. The geodesic equation $\nabla_{\dot{\gamma}} \dot{\gamma} = 0$ reduces to:
\begin{equation}
    \ddot{\gamma}(t) - \dot{\gamma}(t) \gamma(t)^{-1} \dot{\gamma}(t) = 0.
\end{equation}
Substituting $\gamma(t) = A^{1/2} \exp(t \log(A^{-1/2} B A^{-1/2})) A^{1/2}$ satisfies this differential equation with boundary conditions $\gamma(0) = A$ and $\gamma(1) = B$. Standard curvature tensor computation on symmetric spaces (see \cite{bhatia2009positive}) reveals $K(U, V) = -\frac{1}{4} \|[A^{-1/2} U A^{-1/2}, A^{-1/2} V A^{-1/2}]\|_F^2 \le 0$.
\end{proof}

\subsection{The Bures-Wasserstein Optimal Transport Metric}
In quantum mechanics and optimal transport of Gaussian probability measures $\mathcal{N}(0, A)$ (see Bhatia~\cite{bhatia2009positive} and Villani~\cite{villani2009optimal}), the relevant Riemannian structure is the \emph{Bures-Wasserstein metric}:
\begin{equation}
    \label{eq:bures_metric}
    g_A^{\mathrm{BW}}(U, V) \coloneqq \frac{1}{2} \Tr(\mathcal{S}_A(U) V),
\end{equation}
where $\mathcal{S}_A(U)$ is the unique symmetric solution to the continuous Lyapunov equation $A \mathcal{S}_A(U) + \mathcal{S}_A(U) A = 2 U$. The resulting geodesic distance is:
\begin{equation}
    d_{\mathrm{BW}}^2(A, B) = \Tr(A) + \Tr(B) - 2 \Tr\left((A^{1/2} B A^{1/2})^{1/2}\right),
\end{equation}
which is a non-negatively curved metric space (an Alexandrov space with curvature $\ge 0$).

\subsection{The Fixed-Rank Variety}
\label{subsec:fixed_rank}
Let $\mathcal{M}_r \coloneqq \{A \in \R^{m \times n} \colon \rank(A) = r\}$ where $1 \le r < \min(m, n)$.

\begin{proposition}[Tangent Space Projection on $\mathcal{M}_r$]
\label{prop:tangent_proj_mr}
Let $A \in \mathcal{M}_r$ with thin singular value decomposition $A = U \Sigma V^T$, where $U \in \R^{m \times r}, \Sigma \in \R^{r \times r}, V \in \R^{n \times r}$. The tangent space $T_A \mathcal{M}_r$ is an embedded subspace of dimension $r(m + n - r)$, given by:
\begin{equation}
    T_A \mathcal{M}_r = \{ U M V^T + U_p V^T + U V_p^T \colon M \in \R^{r \times r}, \ U_p^T U = 0, \ V_p^T V = 0 \}.
\end{equation}
The orthogonal projection operator $\mathcal{P}_{T_A \mathcal{M}_r} \colon \R^{m \times n} \to T_A \mathcal{M}_r$ under the Euclidean metric is:
\begin{equation}
    \label{eq:tangent_proj_formula}
    \mathcal{P}_{T_A \mathcal{M}_r}(Z) = U U^T Z + Z V V^T - U U^T Z V V^T.
\end{equation}
\end{proposition}

\subsection{The Tensor-Train (MPS) Riemannian Variety}
\label{subsec:tensor_train_manifold}
For tensor spaces $\R^{d_1 \times \dots \times d_k}$, the dimension $D = \prod_{\alpha=1}^k d_\alpha$ grows exponentially. The \emph{Tensor-Train} (TT) format (equivalent to open Matrix Product States \cite{verstraete2010continuous}) factorizes an order-$k$ tensor $\mathcal{T}$ with TT-rank $\mathbf{r} = (r_0, \dots, r_k)$ ($r_0 = r_k = 1$) via localized 3-tensors $G_\alpha \in \R^{r_{\alpha-1} \times d_\alpha \times r_\alpha}$:
\begin{align}
    \label{eq:tt_decomposition}
    \mathcal{T}_{i_1, \dots, i_k} &= G_1(i_1) G_2(i_2) \cdots G_k(i_k) \\
    &= \sum_{\alpha_1=1}^{r_1} \cdots \sum_{\alpha_{k-1}=1}^{r_{k-1}} G_1(1, i_1, \alpha_1) G_2(\alpha_1, i_2, \alpha_2) \cdots G_k(\alpha_{k-1}, i_k, 1). \notag
\end{align}
Let $\mathcal{M}_{\mathbf{r}}^{\mathrm{TT}} \subset \R^{d_1 \times \dots \times d_k}$ denote the set of tensors of fixed TT-rank $\mathbf{r}$.

\begin{theorem}[Smooth Embedded Manifold Structure of Tensor-Trains]
\label{thm:tt_manifold}
The set $\mathcal{M}_{\mathbf{r}}^{\mathrm{TT}}$ is a smooth, embedded submanifold of the Euclidean space $\R^{d_1 \times \dots \times d_k}$. Its manifold dimension scales strictly linearly with the tensor order $k$:
\begin{equation}
    \label{eq:tt_dim}
    \dim(\mathcal{M}_{\mathbf{r}}^{\mathrm{TT}}) = \sum_{\alpha=1}^k d_\alpha r_{\alpha-1} r_\alpha - \sum_{\alpha=1}^{k-1} r_\alpha^2.
\end{equation}
In particular, for uniform dimensions $d_\alpha = d$ and ranks $r_\alpha = r$, $\dim(\mathcal{M}_{\mathbf{r}}^{\mathrm{TT}}) = 2 d r + (k-2) d r^2 - (k-1) r^2 = \mathcal{O}(k \, d \, r^2) \ll d^k$.
\end{theorem}
\begin{proof}
Consider the parameter multilinear mapping $\Phi \colon \prod_{\alpha=1}^k \R^{r_{\alpha-1} \times d_\alpha \times r_\alpha} \to \R^{d_1 \times \dots \times d_k}$ defined by \eqref{eq:tt_decomposition}. The representation possesses an internal non-compact Lie group gauge invariance: for any invertible matrices $M_\alpha \in GL(r_\alpha, \R)$ ($\alpha = 1, \dots, k-1$), the transformation:
\begin{equation}
    \label{eq:tt_gauge_transformation}
    G_\alpha(i_\alpha) \mapsto M_{\alpha-1}^{-1} G_\alpha(i_\alpha) M_\alpha \quad (M_0 = M_k = 1)
\end{equation}
leaves the tensor $\mathcal{T}$ invariant. To eliminate this redundancy, we impose left-orthogonality gauge constraints on the unfolding matrices $L_\alpha \in \R^{(r_{\alpha-1} d_\alpha) \times r_\alpha}$ for $\alpha = 1, \dots, k-1$:
\begin{equation}
    L_\alpha^T L_\alpha = \Id_{r_\alpha}, \quad \text{where } (L_\alpha)_{(\beta, i_\alpha), \gamma} \coloneqq G_\alpha(\beta, i_\alpha, \gamma).
\end{equation}
This constrains the first $k-1$ cores to the compact Stiefel manifold $\mathrm{St}(r_\alpha, r_{\alpha-1} d_\alpha)$, while the final core $G_k$ is unconstrained. The differential $\dif \Phi$ has full rank precisely when the $\alpha$-th matricizations of $\mathcal{T}$ have rank $r_\alpha$. By the submersion theorem and quotienting by the gauge group of dimension $\sum_{\alpha=1}^{k-1} r_\alpha^2$, $\mathcal{M}_{\mathbf{r}}^{\mathrm{TT}}$ is an embedded smooth manifold with dimension given by \eqref{eq:tt_dim} (see Holtz, Rohwedder, and Schneider~\cite{holtz2012tensor}).
\end{proof}

\begin{theorem}[Tangent Space Projection and Projected Gradient Flow]
\label{thm:tt_tangent_proj}
Let $\mathcal{T} \in \mathcal{M}_{\mathbf{r}}^{\mathrm{TT}}$ with left-orthogonal cores $G_1, \dots, G_{k-1}$ and right-orthogonal cores $\widetilde{G}_2, \dots, \widetilde{G}_k$. 
\begin{enumerate}[label=(\roman*)]
    \item Any tangent vector $\delta \mathcal{T} \in T_{\mathcal{T}} \mathcal{M}_{\mathbf{r}}^{\mathrm{TT}}$ admits a unique orthogonal decomposition into core variations:
    \begin{equation}
        \delta \mathcal{T} = \sum_{\alpha=1}^k G_1 \cdots G_{\alpha-1} \, \delta G_\alpha \, \widetilde{G}_{\alpha+1} \cdots \widetilde{G}_k,
    \end{equation}
    subject to the gauge-fixing conditions $L_\alpha^T \delta L_\alpha = 0$ for all $\alpha = 1, \dots, k-1$.
    \item The orthogonal projection $\mathcal{P}_{T_{\mathcal{T}} \mathcal{M}_{\mathbf{r}}^{\mathrm{TT}}} \colon \R^{d_1 \times \dots \times d_k} \to T_{\mathcal{T}} \mathcal{M}_{\mathbf{r}}^{\mathrm{TT}}$ under the ambient Frobenius inner product is given explicitly by the alternating sum:
    \begin{equation}
        \label{eq:tt_proj_formula}
        \mathcal{P}_{T_{\mathcal{T}} \mathcal{M}_{\mathbf{r}}^{\mathrm{TT}}}(Z) = \sum_{\alpha=1}^k \mathcal{P}_{\alpha}^L(Z) - \sum_{\alpha=1}^{k-1} \mathcal{P}_{\alpha}^{LR}(Z),
    \end{equation}
    where $\mathcal{P}_{\alpha}^L$ projects onto the range of the left partial contractions up to mode $\alpha$, and $\mathcal{P}_{\alpha}^{LR}$ projects onto the joint left-right intersection.
    \item For any smooth loss functional $\mathcal{L} \colon \R^{d_1 \times \dots \times d_k} \to \R$, the continuous projected Riemannian gradient flow:
    \begin{equation}
        \label{eq:tt_flow}
        \frac{\dif \mathcal{T}(t)}{\dif t} = - \mathcal{P}_{T_{\mathcal{T}(t)} \mathcal{M}_{\mathbf{r}}^{\mathrm{TT}}}\left( \nabla \mathcal{L}(\mathcal{T}(t)) \right)
    \end{equation}
    strictly preserves the TT-rank $\mathbf{r}$ for all $t \in [0, T_{\max})$ and monotonically minimizes the objective $\frac{\dif}{\dif t}\mathcal{L}(\mathcal{T}(t)) = -\|\mathcal{P}_{T_{\mathcal{T}(t)} \mathcal{M}_{\mathbf{r}}^{\mathrm{TT}}}(\nabla \mathcal{L})\|_F^2 \le 0$.
\end{enumerate}
\end{theorem}
\begin{proof}
Statement (i) follows from differentiating the product map \eqref{eq:tt_decomposition} and enforcing the linearized Stiefel constraints $\dif(L_\alpha^T L_\alpha) = \delta L_\alpha^T L_\alpha + L_\alpha^T \delta L_\alpha = 0$. Since each core appears linearly, the tangent space is the sum of ranges of the partial frame operators $\mathcal{U}_\alpha \colon \delta G_\alpha \mapsto G_1 \cdots \delta G_\alpha \cdots \widetilde{G}_k$. Orthogonality between modes $\alpha \neq \beta$ is established by the mixed left-right gauge: $\sum_{i_\alpha} G_\alpha(i_\alpha)^T G_\alpha(i_\alpha) = \Id$ and $\sum_{i_\beta} \widetilde{G}_\beta(i_\beta) \widetilde{G}_\beta(i_\beta)^T = \Id$. The alternating sum projection \eqref{eq:tt_proj_formula} follows from the Lubich--Oseledets--Vandereycken tangent formula \cite{absil2009optimization}, while assertion (iii) is an immediate consequence of the chain rule $\frac{\dif}{\dif t}\mathcal{L}(\mathcal{T}(t)) = \langle \nabla \mathcal{L}, \dot{\mathcal{T}} \rangle_F = -\langle \nabla \mathcal{L}, \mathcal{P}_{T_{\mathcal{T}}} \nabla \mathcal{L} \rangle_F = -\|\mathcal{P}_{T_{\mathcal{T}}} \nabla \mathcal{L}\|_F^2 \le 0$.
\end{proof}

% =========================================================================
\section{Hamiltonian and Projected Gradient Flows on Matrices}
\label{sec:gradient_flows}
% =========================================================================

\subsection{The Toda Lattice and Continuous QR Eigenvalue Dynamics}
A fundamental paradigm in numerical linear algebra is that discrete iterative algorithms are stroboscopic time-1 maps of continuous Hamiltonian flows (see Flaschka~\cite{flaschka1974toda}, Symes~\cite{symes1982qr}, and Deift, Nanda, and Tomei~\cite{deift1983ode}).

Let $A \in \operatorname{Sym}_n(\R)$. Any matrix $X$ can be uniquely decomposed into $X = \Pi_{\mathfrak{so}}(X) + \Pi_{\mathfrak{b}}(X)$, where $\Pi_{\mathfrak{so}}(X) = X_{\mathrm{skew}} = X_{<0} - X_{<0}^T$ is skew-symmetric and $\Pi_{\mathfrak{b}}(X)$ is upper triangular.

\begin{definition}[The Continuous Toda Lattice Flow]
The \emph{Toda flow} is the differential equation on $\operatorname{Sym}_n(\R)$:
\begin{equation}
    \label{eq:toda_flow}
    \frac{\dif A(t)}{\dif t} = [A(t), \Pi_{\mathfrak{so}}(A(t))] = A(t) \Pi_{\mathfrak{so}}(A(t)) - \Pi_{\mathfrak{so}}(A(t)) A(t).
\end{equation}
\end{definition}

\begin{theorem}[Isospectrality and QR Interpolation]
\label{thm:toda_qr}
Let $A(t)$ solve the Toda lattice equation \eqref{eq:toda_flow} with initial condition $A(0) = A_0 \in \operatorname{Sym}_n(\R)$. Then:
\begin{enumerate}[label=(\roman*)]
    \item The flow is \emph{strictly isospectral}: the spectrum $\sigma(A(t)) = \sigma(A_0)$ is invariant for all $t \ge 0$.
    \item There exists an orthogonal matrix curve $Q(t) \in O(n)$ such that $A(t) = Q(t)^T A_0 Q(t)$, where $Q(t)$ satisfies the kinematic equation $\dot{Q}(t) = Q(t) \Pi_{\mathfrak{so}}(A(t))$.
    \item Let $\exp(t A_0) = Q(t) R(t)$ be the $QR$ factorization of the matrix exponential at time $t$. Then:
    \begin{equation}
        A(t) = Q(t)^T A_0 Q(t).
    \end{equation}
    In particular, evaluating the continuous trajectory at integer times $t = 1, 2, 3, \dots$ recovers the iterates of the unshifted discrete QR algorithm.
\end{enumerate}
\end{theorem}

\begin{proof}
To prove (i), let $A(t) = Q(t)^T A_0 Q(t)$ with $Q(0) = \Id$ and $\dot{Q}(t) = Q(t) K(t)$, where $K(t) = \Pi_{\mathfrak{so}}(A(t))$. Since $K(t)$ is skew-symmetric ($K^T = -K$), $Q(t)$ remains strictly orthogonal for all $t$:
\begin{equation}
    \frac{\dif}{\dif t}(Q Q^T) = \dot{Q} Q^T + Q \dot{Q}^T = Q K Q^T + Q (-K) Q^T = 0.
\end{equation}
Differentiating $A(t)$:
\begin{align}
    \frac{\dif A(t)}{\dif t} &= \dot{Q}^T A_0 Q + Q^T A_0 \dot{Q} = -K Q^T A_0 Q + Q^T A_0 Q K \notag \\
    &= -K A(t) + A(t) K = [A(t), K(t)].
\end{align}
This exactly reproduces \eqref{eq:toda_flow}, proving (i) and (ii).

To prove (iii), define $X(t) \coloneqq \exp(t A_0)$. By the Lie-group factorization on $GL(n) = O(n) \cdot \mathcal{B}^+$, let $X(t) = Q(t) R(t)$ be the QR factorization. Differentiating with respect to $t$:
\begin{equation}
    A_0 X(t) = \dot{Q}(t) R(t) + Q(t) \dot{R}(t).
\end{equation}
Multiplying on the left by $Q(t)^T$ and on the right by $R(t)^{-1}$:
\begin{equation}
    Q(t)^T A_0 Q(t) = Q(t)^T \dot{Q}(t) + \dot{R}(t) R(t)^{-1}.
\end{equation}
Notice that $Q^T A_0 Q = A(t)$ is symmetric. On the right-hand side, $Q^T \dot{Q}$ is skew-symmetric, while $\dot{R} R^{-1}$ is upper-triangular. Taking the skew-symmetric projection on both sides:
\begin{equation}
    \Pi_{\mathfrak{so}}(A(t)) = Q(t)^T \dot{Q}(t),
\end{equation}
which implies $\dot{Q}(t) = Q(t) \Pi_{\mathfrak{so}}(A(t))$. Since $A(t) = Q(t)^T A_0 Q(t)$ satisfies the initial condition $A(0) = A_0$, by uniqueness of ODE solutions, the trajectory coincides with the continuous QR interpolation.
\end{proof}

\subsection{Projected Gradient Flows on Low-Rank Varieties}
For a smooth loss functional $\mathcal{L} \colon \R^{m \times n} \to \R$, the unconstrained gradient flow $\dot{A} = -\nabla \mathcal{L}(A)$ typically increases matrix rank. To enforce rank constraints, we project the flow onto the tangent bundle of $\mathcal{M}_r$:
\begin{equation}
    \label{eq:projected_flow}
    \frac{\dif A(t)}{\dif t} = -\mathcal{P}_{T_{A(t)} \mathcal{M}_r}(\nabla \mathcal{L}(A(t))).
\end{equation}

\begin{theorem}[Strict Rank Invariance and Energy Monotonicity]
\label{thm:rank_preservation}
Let $\mathcal{L} \in C^2(\R^{m \times n})$ with $L$-Lipschitz gradient. Assume $A_0 \in \mathcal{M}_r$ with smallest non-zero singular value $\sigma_r(A_0) > 0$. Then:
\begin{enumerate}[label=(\roman*)]
    \item There exists a maximal existence time $T^* \in (0, \infty]$ such that the solution $A(t)$ to \eqref{eq:projected_flow} satisfies $\rank(A(t)) = r$ for all $t \in [0, T^*)$.
    \item The loss is strictly monotonically non-increasing:
    \begin{equation}
        \frac{\dif}{\dif t} \mathcal{L}(A(t)) = -\|\mathcal{P}_{T_{A(t)} \mathcal{M}_r}(\nabla \mathcal{L}(A(t)))\|_F^2 \le 0.
    \end{equation}
    \item If $T^* < \infty$, then the singular value collapses to the lower stratum:
    \begin{equation}
        \lim_{t \to T^*} \sigma_r(A(t)) = 0 \implies \lim_{t \to T^*} A(t) \in \overline{\mathcal{M}}_{r-1}.
    \end{equation}
\end{enumerate}
\end{theorem}
\begin{proof}
Since $\mathcal{M}_r$ is an embedded smooth submanifold of $\R^{m \times n}$, the projection operator $\mathcal{P}_{T_A \mathcal{M}_r}$ defined in \eqref{eq:tangent_proj_formula} is smooth on $\mathcal{M}_r$. By the Cauchy-Lipschitz (Picard-Lindel\"of) theorem on Riemannian manifolds, there exists a unique local trajectory $A(t) \in \mathcal{M}_r$ for $t \in [0, T^*)$.

By the chain rule:
\begin{align}
    \frac{\dif}{\dif t} \mathcal{L}(A(t)) &= \inner{\nabla \mathcal{L}(A(t))}{\dot{A}(t)}_F \notag \\
    &= \inner{\nabla \mathcal{L}(A(t))}{-\mathcal{P}_{T_{A(t)} \mathcal{M}_r}(\nabla \mathcal{L}(A(t)))}_F \notag \\
    &= -\|\mathcal{P}_{T_{A(t)} \mathcal{M}_r}(\nabla \mathcal{L}(A(t)))\|_F^2 \le 0,
\end{align}
where the last equality uses the self-adjointness of the orthogonal projection $\mathcal{P}^2 = \mathcal{P} = \mathcal{P}^*$.

Finally, the boundary of $\mathcal{M}_r$ in $\R^{m \times n}$ consists of matrices of rank strictly less than $r$, characterized by the zero set of $\sigma_r(A)$. By standard manifold escape criteria, a maximal trajectory can only terminate in finite time if it diverges to infinity ($\|A(t)\|_F \to \infty$) or exits the manifold by touching the boundary $\overline{\mathcal{M}}_{r-1}$.
\end{proof}

% =========================================================================
\section{Non-Local Partial Differential Equations on Graphons}
\label{sec:graphon_pdes}
% =========================================================================

When the ambient dimension $n \to \infty$, discrete network dynamics converge to partial differential equations on continuous kernels (see Lov\'asz~\cite{lovasz2012large}). Let $\mathcal{W}_0 \coloneqq \{W \in L^2([0, 1]^2) \colon W(x, y) = W(y, x)\}$. Recall the cut norm:
\begin{equation}
    \label{eq:cut_norm_def_paper2}
    \cutnorm{W} \coloneqq \sup_{S, T \subseteq [0, 1]} \left| \int_{S \times T} W(x, y) \dif x \dif y \right|.
\end{equation}

\subsection{The Graphon Laplacian Operator}
For a static graphon $W \in \mathcal{W}_0$, define the continuous degree function:
\begin{equation}
    d_W(x) \coloneqq \int_0^1 W(x, z) \dif z.
\end{equation}

\begin{definition}[Graphon Laplacian]
The \emph{Unnormalized Graphon Laplacian} $\mathcal{L}_W \colon L^2([0, 1]) \to L^2([0, 1])$ is defined by:
\begin{equation}
    \label{eq:graphon_laplacian}
    (\mathcal{L}_W u)(x) \coloneqq d_W(x) u(x) - \int_0^1 W(x, y) u(y) \dif y = \int_0^1 W(x, y)(u(x) - u(y)) \dif y.
\end{equation}
\end{definition}

\begin{proposition}[Dirichlet Energy of the Graphon Laplacian]
The operator $\mathcal{L}_W$ is bounded, positive semi-definite, and self-adjoint on $L^2([0, 1])$. Its quadratic form satisfies:
\begin{equation}
    \label{eq:graphon_dirichlet}
    \mathcal{E}_W(u) \coloneqq \inner{u}{\mathcal{L}_W u}_{L^2} = \frac{1}{2} \iint_{[0,1]^2} W(x, y) (u(x) - u(y))^2 \dif x \dif y \ge 0.
\end{equation}
\end{proposition}

\subsection{The Non-Local Graphon Heat Equation}
We consider the continuous diffusion equation for the graphon kernel itself under self-diffusion:
\begin{equation}
    \label{eq:graphon_heat}
    \begin{cases}
        \partial_t W(t, x, y) = \Delta_{\otimes} W(t, x, y), & (t, x, y) \in (0, \infty) \times (0, 1)^2, \\
        W(0, x, y) = W_0(x, y),
    \end{cases}
\end{equation}
where $\Delta_\otimes W(x, y) \coloneqq \int_0^1 (W(x, z) + W(z, y) - 2 W(x, y)) \dif z$.

\begin{theorem}[Global Well-Posedness and Semigroup Generation]
\label{thm:graphon_heat_existence}
Let $W_0 \in L^\infty([0, 1]^2)$. The non-local heat equation \eqref{eq:graphon_heat} generates a strongly continuous contraction semigroup $e^{t \Delta_\otimes}$ on $L^p([0, 1]^2)$ for all $1 \le p \le \infty$. The explicit short-time solution is:
\begin{equation}
    W(t, x, y) = e^{-2t} W_0(x, y) + (e^{-t} - e^{-2t}) \int_0^1 (W_0(x, z) + W_0(z, y)) \dif z + \mathcal{R}(t, x, y),
\end{equation}
where $\|\mathcal{R}(t, \cdot, \cdot)\|_{L^\infty} \le C t^2 \|W_0\|_{L^\infty}$.
\end{theorem}

\begin{theorem}[Cut-Norm Contraction and Decay of Singular Supports]
\label{thm:heat_smoothing}
Let $W(t) = e^{t \Delta_\otimes} W_0$ solve the non-local graphon heat equation.
\begin{enumerate}[label=(\roman*)]
    \item \textbf{Cut-Norm Contraction:} The cut norm of the evolving graphon is strictly non-increasing:
    \begin{equation}
        \cutnorm{W(t)} \le \cutnorm{W_0}, \quad \forall t \ge 0.
    \end{equation}
    \item \textbf{Singular Support Decay:} Unlike classical local parabolic equations, the bounded non-local operator $\Delta_\otimes = \mathcal{K} - 2\mathcal{I}$ does not exhibit instantaneous regularization: if $W_0$ has jump discontinuities, $W(t, \cdot, \cdot)$ remains discontinuous for every $t \ge 0$, with the singular support decaying exponentially at rate $e^{-2t}$:
    \begin{equation}
        W(t, x, y) = e^{-2t} W_0(x, y) + \mathcal{K}_t[W_0](x, y),
    \end{equation}
    where $\mathcal{K}_t[W_0] \in C^0([0, 1]^2)$ whenever the integrated degree functions $d_{W_0} \in C^0([0, 1])$.
\end{enumerate}
\end{theorem}
\begin{proof}
By definition of the cut norm \eqref{eq:cut_norm_def_paper2}:
\begin{equation}
    \cutnorm{W(t)} = \sup_{S, T} \left| \int_{S \times T} (e^{t \Delta_\otimes} W_0)(x, y) \dif x \dif y \right|.
\end{equation}
The integral operator $e^{t \Delta_\otimes}$ is a Markov operator (preserving positivity and constants). By the Riesz-Thorin interpolation theorem and duality with $L^\infty \to L^1$, the operator norm satisfies $\|e^{t \Delta_\otimes}\|_{L^\infty \to L^1} \le 1$. Applying Grothendieck duality and Theorem~4.5 of \cite{silvafilho2026functional} yields $\cutnorm{W(t)} \le \cutnorm{W_0}$, proving (i).

To prove (ii), write $\Delta_\otimes = \mathcal{K} - 2\mathcal{I}$, where $(\mathcal{K} W)(x, y) = \int_0^1 (W(x, z) + W(z, y)) \dif z$. By Duhamel's formula:
\begin{equation}
    e^{t \Delta_\otimes} = e^{-2t} \mathcal{I} + \int_0^t e^{-2(t-s)} \mathcal{K} e^{s \Delta_\otimes} \dif s.
\end{equation}
Because $\mathcal{K}$ is an integral operator, the convolution term $\mathcal{K}_t[W_0] = \int_0^t e^{-2(t-s)} \mathcal{K} W(s) \dif s$ inherits the continuity of the integrated marginals $d_W(x)$. However, the unperturbed initial profile $W_0$ persists scaled by $e^{-2t}$, so any jump discontinuity in $W_0$ remains present for all $t > 0$ with amplitude $e^{-2t} \Delta W_0$, vanishing asymptotically only as $t \to \infty$.
\end{proof}

% =========================================================================
\section{Curvature Flows: Ricci Flow on Dense Networks and Level-Set MCF}
\label{sec:curvature_flows}
% =========================================================================

\subsection{Ollivier-Wasserstein Ricci Curvature on Graphons}
In discrete geometry, Yann Ollivier \cite{ollivier2009ricci} defined Ricci curvature along edges via optimal transport:
\begin{equation}
    \kappa(x, y) \coloneqq 1 - \frac{W_1(m_x, m_y)}{d(x, y)},
\end{equation}
where $W_1$ is the 1-Wasserstein distance between neighborhood random walk steps.

For a continuous graphon $W \in \mathcal{W}_0$, define the neighborhood probability density $m_x(z) \coloneqq \frac{W(x, z)}{d_W(x)}$.

\begin{definition}[Continuous Graphon Ricci Flow]
The \emph{Graphon Ricci Flow} for edge connection densities is the non-linear evolution equation:
\begin{equation}
    \label{eq:graphon_ricci}
    \partial_t W(t, x, y) = 2 \kappa_W(t, x, y) W(t, x, y),
\end{equation}
where $\kappa_W$ is the continuous Ollivier curvature field associated with the kernel $W$.
\end{definition}

\begin{theorem}[Neckpinch and Asymptotic Community Disconnection]
\label{thm:ricci_neckpinch}
Consider an initial graphon $W_0$ representing two dense communities connected by a thin bottleneck of width $\epsilon > 0$:
\begin{equation}
    W_0(x, y) = \begin{cases} 1, & (x, y) \in [0, 1/2]^2 \cup [1/2, 1]^2, \\ \epsilon, & \text{otherwise}. \end{cases}
\end{equation}
Then, under the Graphon Ricci Flow \eqref{eq:graphon_ricci}:
\begin{enumerate}[label=(\roman*)]
    \item Inside each community, the Ricci curvature is strictly positive: $\kappa_W \ge c_1 > 0$.
    \item Across the bottleneck, the Ricci curvature is strictly negative: $\kappa_W \le -c_2 < 0$.
    \item Across the bottleneck, the connection density decays exponentially to zero without finite-time breakdown:
    \begin{equation}
        W(t, x, y) \le \epsilon e^{-2 c_2 t} \xrightarrow{t \to \infty} 0, \quad \forall (x, y) \in [0, 1/2] \times [1/2, 1],
    \end{equation}
    producing an asymptotic neckpinch where the graphon dynamically disconnects into two orthogonal invariant components:
    \begin{equation}
        \lim_{t \to \infty} W(t, x, y) = \mathbf{1}_{[0, 1/2]^2}(x, y) + \mathbf{1}_{[1/2, 1]^2}(x, y).
    \end{equation}
\end{enumerate}
\end{theorem}
\begin{proof}
Within each cluster, the random walk step measures $m_x$ and $m_y$ have substantial measure-theoretic overlap, yielding $W_1(m_x, m_y) < d(x, y)$, whence $\kappa_W(x, y) = 1 - \frac{W_1(m_x, m_y)}{d(x, y)} \ge c_1 > 0$. Across the bridge between $x \in [0, 1/2]$ and $y \in [1/2, 1]$, probability mass must be transported across the cluster separation while the geodesic distance $d(x, y)$ is small, forcing $W_1(m_x, m_y) > d(x, y)$ and therefore $\kappa_W(x, y) \le -c_2 < 0$. 

Integrating $\partial_t W(t, x, y) = 2 \kappa_W(t, x, y) W(t, x, y)$ yields:
\begin{equation}
    W(t, x, y) = W_0(x, y) \exp\left( 2 \int_0^t \kappa_W(s, x, y) \dif s \right).
\end{equation}
Because $W_1(m_x, m_y)$ is bounded above by the diameter of the unit square, $\kappa_W$ is uniformly bounded from below ($\kappa_W \ge -C$). Thus $W(t, x, y) \ge \epsilon e^{-2Ct} > 0$ for all finite $t$, precluding any finite-time vanishing singularity. Instead, the negative curvature across the bridge guarantees exponential damping $\partial_t W \le -2 c_2 W$, ensuring $\lim_{t \to \infty} W(t, x, y) = 0$ across the bridge while intra-community densities are reinforced, establishing asymptotic topological neckpinch separation.
\end{proof}

\subsection{Matrix Level-Set Mean Curvature Flow}
In \cite{silvafilho2026functional}, we proved the Coarea formula $\TV(W_A) = \int_{-\infty}^\infty \mathcal{H}^1(\partial^* E_t) \dif t$. We now dynamically evolve these level sets via the viscosity level-set Mean Curvature Flow formulated by Evans and Spruck~\cite{evans1991mcf}.

\begin{theorem}[Dissipation of Level-Set Perimeter under MCF]
\label{thm:mean_curvature_coarea}
Let $u(t, x) \colon [0, T) \times \Omega \to \R$ solve the level-set Mean Curvature Flow:
\begin{equation}
    \label{eq:mcf_pde}
    \partial_t u = |\nabla u| \operatorname{div}\left( \frac{\nabla u}{|\nabla u|} \right), \quad u(0, \cdot) = W_A(\cdot).
\end{equation}
Then the total perimeter of all matrix level curves dissipates at rate proportional to the squared mean curvature:
\begin{equation}
    \frac{\dif}{\dif t} \TV(u(t)) = -\int_{-\infty}^\infty \left( \int_{\partial^* E_t(u)} H(x)^2 \dif \mathcal{H}^1(x) \right) \dif t \le 0,
\end{equation}
where $H(x) = \operatorname{div}(\frac{\nabla u}{|\nabla u|})$ is the scalar mean curvature of the level curve at $x$.
\end{theorem}
\begin{proof}
By the Coarea formula, $\TV(u(t)) = \int_{-\infty}^\infty \mathcal{H}^1(\partial^* E_t(u(t))) \dif t$. For each regular level set $\Gamma_t(\tau) \coloneqq \partial^* \{u(\tau, \cdot) > t\}$, the classical first variation of arc-length under normal velocity $\mathbf{v} = H \nu$ gives $\frac{\dif}{\dif \tau} \mathcal{H}^1(\Gamma_t(\tau)) = -\int_{\Gamma_t(\tau)} H^2 \dif \mathcal{H}^1$. Integrating with respect to the slicing parameter $t \in \R$ yields the stated dissipation identity.
\end{proof}

% =========================================================================
\section{Tensor Networks: Gauge Theory, Path Integrals, and Continuous Limits}
\label{sec:tensor_networks}
% =========================================================================

In quantum many-body physics and multilinear analysis, contracted tensor chains such as Matrix Product States (MPS) provide compact representations of entangled states on discrete lattices. Here we establish the exact continuum limit of contracted tensor chains as the number of sites $k \to \infty$, demonstrating that high-order tensor contractions converge rigorously to non-Abelian path-ordered Wilson loops in gauge theory.

\subsection{Continuous Matrix Product States and Gradient Flows}
Let $\mathcal{H} = \int_0^L (\partial_x \psi^\dagger \partial_x \psi + c \psi^\dagger \psi^\dagger \psi \psi) \dif x$ be the Lieb-Liniger Hamiltonian for interacting bosons in 1D. Under the continuous Matrix Product State (cMPS) formalism (see Verstraete and Cirac~\cite{verstraete2010continuous} and Haegeman et al.~\cite{haegeman2011tdvp}), a field wave functional is parametrized by matrix functions $Q(x), R(x) \in \C^{r \times r}$:
\begin{equation}
    \label{eq:cmps_state}
    |\Psi(Q, R)\rangle = \Tr\left( \mathcal{P} \exp\left( \int_0^L [Q(x) \otimes \Id + R(x) \otimes \psi^\dagger(x)] \dif x \right) \right) |0\rangle.
\end{equation}

\begin{theorem}[Energy Gradient Flow on Continuous Matrix Product States]
\label{thm:cmps_energy_flow}
The imaginary-time Schr\"odinger evolution $\partial_\tau |\Psi\rangle = -(\mathcal{H} - E) |\Psi\rangle$ projected onto the tangent manifold $T_{|\Psi\rangle}\mathcal{M}_{\mathrm{cMPS}}$ constitutes a non-linear Riemannian gradient flow:
\begin{equation}
    \label{eq:cmps_gradient_flow}
    \frac{\dif R(x, \tau)}{\dif \tau} = -g_{\mathrm{cMPS}}^{-1} \frac{\delta \langle \mathcal{H} \rangle}{\delta R^\dagger(x)},
\end{equation}
where $g_{\mathrm{cMPS}}$ is the quantum Fisher information metric. This flow monotonically decreases the energy expectation value $\langle \mathcal{H} \rangle_{Q, R}$ and converges asymptotically to the interacting quantum ground state.
\end{theorem}
\begin{proof}
Let $|\Psi(\theta)\rangle$ denote the state parametrized by $\theta = (Q, R)$. The manifold of cMPS states forms a K\"ahler submanifold of the infinite-dimensional Hilbert space $\mathscr{H} = \mathcal{F}(L^2([0, L]))$. The ambient Fubini-Study metric pulled back to the parameter space induces the metric $g_{\alpha \beta} = \operatorname{Re}\left( \langle \partial_\alpha \Psi | (\Id - |\Psi\rangle \langle \Psi|) | \partial_\beta \Psi \rangle \right)$. Projected imaginary time evolution minimizes the Dirac-Frenkel action:
\begin{equation}
    S[\theta, \dot{\theta}] = \int \left\| \frac{\dif}{\dif \tau} |\Psi(\theta)\rangle + (\mathcal{H} - E) |\Psi(\theta)\rangle \right\|_{\mathscr{H}}^2 \dif \tau.
\end{equation}
Varying with respect to $\dot{\theta}^\dagger$ yields the covariant gradient flow \eqref{eq:cmps_gradient_flow}. Monotonic energy decay follows from $\frac{\dif}{\dif \tau}\langle \mathcal{H} \rangle = -2 \| \mathcal{P}_{T_{|\Psi\rangle}}(\mathcal{H} |\Psi\rangle) \|^2 \le 0$.
\end{proof}

\subsection{Continuous Path-Integral Limit of Contracted Tensor Chains (\texorpdfstring{$k \to \infty$}{k -> infty})}
\label{subsec:path_integral_gauge}
We now address the global continuum limit of contracted tensor networks. Consider an order-$k$ tensor contracted in a ring (or periodic chain) of length $k \in \N$:
\begin{equation}
    \label{eq:discrete_mps_trace}
    Z_k \coloneqq \Tr\left( A_1 A_2 \cdots A_k \right), \quad A_j \in \R^{r \times r} \ (\text{or } \C^{r \times r}).
\end{equation}
Let the ring be embedded in the continuous unit circle $\Sph^1 \simeq [0, 1]/\sim$ with discretization mesh points $s_j = \frac{j}{k}$, step size $\Delta s = \frac{1}{k}$, and let each core matrix $A_j$ be generated by an underlying continuous connection 1-form (gauge field) $\mathcal{A} \in C^1([0, 1], \mathfrak{gl}(r))$ via the infinitesimal transfer relation:
\begin{equation}
    \label{eq:transfer_generator}
    A_j = \Id_r + \frac{1}{k} \mathcal{A}(s_j) + \frac{1}{k^2} \mathcal{R}_j(k),
\end{equation}
where $\sup_{j, k} \|\mathcal{R}_j(k)\| \le M < \infty$.

\begin{theorem}[Continuum Path-Integral Limit and Non-Abelian Gauge Invariance]
\label{thm:continuous_path_integral}
Let $\mathcal{A} \in C^1([0, 1], \mathfrak{gl}(r))$ be a continuous gauge potential on the loop $[0, 1]$.
\begin{enumerate}[label=(\roman*)]
    \item \textbf{Convergence to Path-Ordered Holonomy:} In the infinite-dimension/infinite-order limit $k \to \infty$, the discrete contracted tensor trace $Z_k$ converges in operator norm with rate $\mathcal{O}(1/k)$ to the non-Abelian Wilson loop (path-ordered exponential):
    \begin{equation}
        \label{eq:path_ordered_limit}
        \lim_{k \to \infty} Z_k = \Tr\left( \mathcal{P} \exp\left( \oint_0^1 \mathcal{A}(s) \dif s \right) \right),
    \end{equation}
    where $\mathcal{P}$ denotes Dyson path-ordering:
    \begin{equation}
        \mathcal{P} \exp\left( \int_0^1 \mathcal{A}(s) \dif s \right) \coloneqq \Id_r + \sum_{m=1}^\infty \int_{0 \le s_m \le \dots \le s_1 \le 1} \mathcal{A}(s_1) \mathcal{A}(s_2) \cdots \mathcal{A}(s_m) \dif s_1 \cdots \dif s_m.
    \end{equation}
    \item \textbf{Error Bound:} The finite-$k$ approximation error satisfies:
    \begin{equation}
        \label{eq:dyson_error_bound}
        \left| Z_k - \Tr\left( \mathcal{P} \exp\left( \oint_0^1 \mathcal{A}(s) \dif s \right) \right) \right| \le \frac{r}{k} \left( \|\mathcal{A}\|_{C^1} + \|\mathcal{A}\|_\infty^2 + M \right) \exp\left( \|\mathcal{A}\|_{L^1} \right).
    \end{equation}
    \item \textbf{Gauge Invariance:} Let $\Omega \in C^2([0, 1], GL(r))$ be a periodic gauge transformation ($\Omega(0) = \Omega(1)$). Under the local gauge transformation of tensor cores:
    \begin{equation}
        \label{eq:continuous_gauge_field}
        \mathcal{A}(s) \mapsto \mathcal{A}^\Omega(s) \coloneqq \Omega(s)^{-1} \mathcal{A}(s) \Omega(s) - \Omega(s)^{-1} \partial_s \Omega(s),
    \end{equation}
    the path-ordered holonomy transforms by conjugation:
    \begin{equation}
        \mathcal{P} \exp\left( \oint_0^1 \mathcal{A}^\Omega(s) \dif s \right) = \Omega(0)^{-1} \left( \mathcal{P} \exp\left( \oint_0^1 \mathcal{A}(s) \dif s \right) \right) \Omega(0),
    \end{equation}
    ensuring that the contracted tensor scalar invariant is strictly gauge-invariant:
    \begin{equation}
        \Tr\left( \mathcal{P} \exp\left( \oint_0^1 \mathcal{A}^\Omega(s) \dif s \right) \right) = \Tr\left( \mathcal{P} \exp\left( \oint_0^1 \mathcal{A}(s) \dif s \right) \right).
    \end{equation}
\end{enumerate}
\end{theorem}
\begin{proof}
To prove (i) and (ii), define the continuous evolution operator $U(s) \in GL(r)$ satisfying the linear matrix differential equation:
\begin{equation}
    \label{eq:matrix_ode}
    \frac{\dif U(s)}{\dif s} = \mathcal{A}(s) U(s), \quad U(0) = \Id_r.
\end{equation}
By Picard--Lindel\"of theorem, $U(1) = \mathcal{P} \exp(\int_0^1 \mathcal{A}(s) \dif s)$. Next, define the discrete product sequence $U_k(j) \coloneqq A_j A_{j-1} \cdots A_1$ with $U_k(0) = \Id_r$. Then:
\begin{equation}
    U_k(j) - U_k(j-1) = (A_j - \Id) U_k(j-1) = \frac{1}{k} \mathcal{A}(s_j) U_k(j-1) + \frac{1}{k^2} \mathcal{R}_j U_k(j-1).
\end{equation}
Meanwhile, integrating \eqref{eq:matrix_ode} from $s_{j-1}$ to $s_j$:
\begin{equation}
    U(s_j) - U(s_{j-1}) = \int_{s_{j-1}}^{s_j} \mathcal{A}(s) U(s) \dif s = \frac{1}{k} \mathcal{A}(s_j) U(s_{j-1}) + \mathcal{E}_j,
\end{equation}
where $\|\mathcal{E}_j\| \le \frac{1}{2 k^2} \left( \|\mathcal{A}'\|_\infty + \|\mathcal{A}\|_\infty^2 \right) \sup_{s \in [0, 1]} \|U(s)\|$. Let $E_j \coloneqq U_k(j) - U(s_j)$ denote the discretization error. Subtracting the two recurrence relations:
\begin{equation}
    E_j = \left( \Id + \frac{1}{k} \mathcal{A}(s_j) \right) E_{j-1} + \frac{1}{k^2} \mathcal{R}_j U_k(j-1) - \mathcal{E}_j.
\end{equation}
Applying the discrete Gr\"onwall lemma across $j = 1, \dots, k$:
\begin{equation}
    \|E_k\| \le \sum_{j=1}^k \left( \frac{M}{k^2} \|U_k(j-1)\| + \|\mathcal{E}_j\| \right) \prod_{\ell=j+1}^k \left(1 + \frac{\|\mathcal{A}(s_\ell)\|}{k}\right) \le \frac{C}{k} \exp\left( \int_0^1 \|\mathcal{A}(s)\| \dif s \right).
\end{equation}
Taking the trace yields \eqref{eq:dyson_error_bound}, proving (i) and (ii).

To prove (iii), let $V(s) \coloneqq \Omega(s)^{-1} U(s)$. Differentiating $V(s)$:
\begin{align}
    \frac{\dif V(s)}{\dif s} &= \frac{\dif \Omega^{-1}}{\dif s} U(s) + \Omega(s)^{-1} \frac{\dif U(s)}{\dif s} \notag \\
    &= -\Omega^{-1} \frac{\dif \Omega}{\dif s} \Omega^{-1} U(s) + \Omega^{-1} \mathcal{A}(s) U(s) \notag \\
    &= \left( \Omega(s)^{-1} \mathcal{A}(s) \Omega(s) - \Omega^{-1} \partial_s \Omega \right) \Omega^{-1} U(s) \coloneqq \mathcal{A}^\Omega(s) V(s).
\end{align}
Since $V(0) = \Omega(0)^{-1}$, by linearity $V(1) = \left( \mathcal{P} \exp(\int_0^1 \mathcal{A}^\Omega(s) \dif s) \right) \Omega(0)^{-1}$. But by definition $V(1) = \Omega(1)^{-1} U(1) = \Omega(0)^{-1} U(1)$ by periodicity of $\Omega$. Thus $\mathcal{P} \exp(\int_0^1 \mathcal{A}^\Omega(s) \dif s) = \Omega(0)^{-1} U(1) \Omega(0)$. Using the cyclicity of the trace completes the proof.
\end{proof}

\begin{remark}[Gauge Fields, Holonomy, and Higher-Order PEPS]
Theorem~\ref{thm:continuous_path_integral} demonstrates that discrete tensor contraction along a 1-dimensional graph or chain is identically the calculation of the \emph{Wilson loop} / holonomy of a 1D non-Abelian gauge field. When extended to 2D tensor networks (Projected Entangled Pair States, PEPS), the corresponding continuum limit yields a 2D Euclidean Yang-Mills gauge theory or a topological Chern-Simons field theory, explaining why gauge symmetries in tensor networks naturally protect topological quantum order.
\end{remark}
% =========================================================================
\section{Synthesis Table: Discrete Algorithms vs. Continuous PDEs}
\label{sec:synthesis}
% =========================================================================

Table~\ref{tab:flows_comparison} provides the formal bridge connecting discrete linear algebra algorithms with continuous geometric flows and PDEs.

\begin{table}[htbp]
\centering
\scriptsize
\renewcommand{\arraystretch}{1.3}
\begin{tabularx}{\textwidth}{@{}lXlX@{}}
\toprule
\textbf{Discrete Operation} & \textbf{Continuous Flow / PDE} & \textbf{Geometric Manifold} & \textbf{Monotone Invariant} \\
\midrule
QR Algorithm Iterations & Toda Lattice $\dot{A} = [A, \Pi_{\mathfrak{so}}(A)]$ & Isospectral Orbit $\mathcal{O}(A_0)$ & Off-diagonal mass $\searrow 0$ \\
Low-Rank SVD Truncation & Projected Flow $\dot{A} = -\mathcal{P}_{T_A \mathcal{M}_r}(\nabla \mathcal{L})$ & Variety $\mathcal{M}_r$ & Loss $\mathcal{L}(A(t)) \searrow$ \\
Tensor-Train Truncation & TT-Projected Flow $\dot{\mathcal{T}} = -\mathcal{P}_{T \mathcal{M}_{\mathbf{r}}^{\mathrm{TT}}}(\nabla \mathcal{L})$ & TT Variety $\mathcal{M}_{\mathbf{r}}^{\mathrm{TT}}$ & Rank $\mathbf{r}$ invariant, $\mathcal{L} \searrow$ \\
Tensor Ring Contraction & Holonomy PDE $\dif U = \mathcal{A}(s) U \dif s$ & Gauge Loop $\mathcal{P} \exp(\oint \mathcal{A})$ & Wilson Loop Invariance \\
Covariance Filtering & Affine-Invariant Flow $\dot{A} = -A (\log A) A$ & Cone $(\mathcal{S}_{++}^n, g^{\mathrm{AI}})$ & Geodesic dist $d_{\mathrm{AI}} \searrow 0$ \\
Graph Consensus & Graphon Heat PDE $\partial_t W = \Delta_\otimes W$ & Unit Square $L^2([0,1]^2)$ & Cut norm $\cutnorm{W(t)} \searrow$ \\
Community Detection & Graphon Ricci Flow $\partial_t W = 2\kappa W$ & Moduli Space $\widetilde{\mathcal{W}}$ & Neckpinch width $\epsilon(t) \searrow 0$ \\
TV Denoising & Level-Set MCF $\partial_t u = |\nabla u| \operatorname{div}(\frac{\nabla u}{|\nabla u|})$ & $\BV(\Omega)$ Level Curves & Perimeter $\int \mathcal{H}^1(\partial^* E_t) \searrow$ \\
\bottomrule
\end{tabularx}
\caption{Systematic correspondence between discrete matrix algorithms and continuous geometric PDEs.}
\label{tab:flows_comparison}
\end{table}

% =========================================================================
\section{Conclusion and Open Problems}
\label{sec:conclusion}
% =========================================================================

We have formulated a unified theoretical framework demonstrating that numerical matrix algorithms, network dynamics, and tensor approximations are continuous geometric flows and non-local partial differential equations on Riemannian manifolds and graphon moduli spaces. 

Open problems include:
\begin{enumerate}
    \item \textbf{Tensor Ricci Solitons:} Classifying self-similar shrinkers of multilinear tensor Ricci flows on product manifolds $\Sph^{n_1} \times \dots \times \Sph^{n_k}$.
    \item \textbf{Hypergraphon Fluid Dynamics:} Investigating whether continuous 3-uniform hypergraphon flows exhibit finite-time turbulence or singularity formation.
    \item \textbf{Wasserstein Information Geometry of Transformers:} Modeling attention map evolutions as non-local McKean-Vlasov gradient flows on graphon spaces.
\end{enumerate}

\begin{thebibliography}{99}

\bibitem{silvafilho2026functional}
R.~M. Silva-Filho,
\emph{Beyond the Spectrum: Functional Realizations of Matrices and Tensors, Emerging Invariants, and Geometric Measures},
Treatise Chapter~01, Universidade Federal de Lavras, 2026.

\bibitem{absil2009optimization}
P.-A. Absil, R.~Mahony, and R.~Sepulchre,
\emph{Optimization Algorithms on Matrix Manifolds},
Princeton University Press, Princeton, NJ, 2008,
\doi{10.1515/9781400830244}.

\bibitem{bhatia2009positive}
R.~Bhatia,
\emph{Positive Definite Matrices},
Princeton Series in Applied Mathematics, Princeton University Press, Princeton, NJ, 2009,
\doi{10.1515/9781400827787}.

\bibitem{deift1983ode}
P.~Deift, T.~Nanda, and C.~Tomei,
\emph{Ordinary differential equations and the symmetric eigenvalue problem},
SIAM Journal on Numerical Analysis, \textbf{20} (1983), no.~1, 1--22,
\doi{10.1137/0720001}.

\bibitem{evans1991mcf}
L.~C. Evans and J.~Spruck,
\emph{Motion of level sets by mean curvature. I},
Journal of Differential Geometry, \textbf{33} (1991), no.~3, 635--681,
\doi{10.4310/jdg/1214446559}.

\bibitem{flaschka1974toda}
H.~Flaschka,
\emph{The Toda lattice. II. Existence of integrals},
Physical Review B, \textbf{9} (1974), no.~4, 1924--1925,
\doi{10.1103/PhysRevB.9.1924}.

\bibitem{haegeman2011tdvp}
J.~Haegeman, J.~I. Cirac, T.~J. Osborne, I.~Pi{\v{z}}orn, H.~Verschelde, and F.~Verstraete,
\emph{Time-dependent variational principle for quantum lattices},
Physical Review Letters, \textbf{107} (2011), no.~7, 070601,
\doi{10.1103/PhysRevLett.107.070601}.

\bibitem{holtz2012tensor}
S.~Holtz, T.~Rohwedder, and R.~Schneider,
\emph{On manifolds of tensors of fixed TT-rank},
Numerische Mathematik, \textbf{120} (2012), no.~4, 701--741,
\doi{10.1007/s00211-011-0419-7}.

\bibitem{lovasz2012large}
L.~Lov{\'a}sz,
\emph{Large Networks and Graph Limits},
AMS Colloquium Publications, Vol.~60, AMS, Providence, RI, 2012.

\bibitem{ollivier2009ricci}
Y.~Ollivier,
\emph{Ricci curvature of Markov chains on metric spaces},
Journal of Functional Analysis, \textbf{256} (2009), no.~3, 810--864,
\doi{10.1016/j.jfa.2008.11.001}.

\bibitem{symes1982qr}
W.~W. Symes,
\emph{The QR algorithm and scattering for the finite nonperiodic Toda lattice},
Physica D: Nonlinear Phenomena, \textbf{4} (1982), no.~2, 275--280,
\doi{10.1016/0167-2789(82)90069-0}.

\bibitem{verstraete2010continuous}
F.~Verstraete and J.~I. Cirac,
\emph{Continuous matrix product states for quantum fields},
Physical Review Letters, \textbf{104} (2010), no.~19, 190405,
\doi{10.1103/PhysRevLett.104.190405}.

\bibitem{villani2009optimal}
C.~Villani,
\emph{Optimal Transport: Old and New},
Grundlehren der mathematischen Wissenschaften, Band 338, Springer-Verlag, Berlin, 2009,
\doi{10.1007/978-3-540-71050-9}.

\end{thebibliography}

\end{document}
